\section{Conclusion and Outlook}

The central question behind this project work is to find out whether a differentiable physics simulator, 
used as the prior mean in bayesian optimisation (BO), can actually help with a realistic accelerator tuning problem (ARES Experimental Area), 
or does it only work on a toy example (FODO lattice)? Based on everything presented, including implementation, experiments and the results, it can confidently be said "yes".
The problem tackled in this project is the transverse beam tuning task
at ARES Experimental Area, the idea of Cheetah-based prior demonstrated on a simple two-dimensional FODO lattice~\cite{kaiser2024bridging} 
is scaled up to a five-dimensional beam tuning task under this project scope, which includes 5 tunable magnets (three quadrupoles and two correctors), six unknown quadrupole
misalignment values that the prior had to learn during optimisation and an objective that asks the
optimiser to focus and centre the beam simultaneously. 
The code is built around a class \texttt{AresPriorMean}, that
wraps Cheetah and plugs it into the Xopt framework. The approach is tested against three experimental conditions. The matched scenario, where
the prior already knows the correct and ideal system state, the mismatched case, where wrong prior is given on purpose, 
forcing it to learn  the correct values during optimisation, and the matched\_prior\_newtask
scenario, where all settings (beam and  quadrupole misalignment values) are same as mismatched case, but the correct prior is given upfront. For experiments conducted under this project work, the third scenario is considered the actual matched case.
Both Cheetah-based BO were compared against two baseline techniques, the standard BO with zero prior and Nelder-Mead Simplex. 

The results confirm that using physics-informed prior in BO, significantly improves the performance. Both the BO prior variants (matched and mismatched) achieved a median Best MAE of \SI{5.75}{\mu}m and \SI{6.71}{\mu}m respectively, which is significantly better than standard BO (45.88$\mu$m) and Nelder-Mead(67.13$\mu$m).
BO with both matched and mismatched prior also achieved better Final MAE compared to baseline methods, with a median of \SI{42.60}{\mu}m and \SI{40.54}{\mu}m respectively, while standard BO and Nelder-Mead had a median Final MAE of \SI{147.96}{\mu}m and \SI{67.84}{\mu}m respectively.
The matched prior hit the 40$\mu$m target in a median of 4 steps, the mismatched prior in 10.5, whereas standard BO required 82.5 steps on rare occasions (20\% success rate) and 
Nelder-Mead never reached the target. What makes the approach especially promising for real world use is its robustness. Even when starting with incorrect misalignment values 
and modified beam distribution, the prior still outperformed the standard BO and reached on a comparable level as the matched prior scenario 
of BO very soon within 15-20 iterations. The trainable misalignment parameters successfully converged very close to the ground truth purely from optimisation data, which means that the method does 
not just tune the beam but it also learns something about the machine's physical state along the way.  \

Ofcourse, the ARES Experimental Area is still a relatively small problem compared to large-scale facilities like European XFEL, 
but this is a significant step forward from the two-dimensional FODO lattice that is previously tested to evaluate cheetah's potential. The fact that the method scales gracefully 
to this level sends quite good signs for scaling further and handling more complex beamlines in future. 



